\section{Extended Methodological Clarifications}

\subsection{What Is the Goal of the Study?}

The goal of the study is to design and evaluate a cooperative caching method for IoT edge networks whose cache decisions account for mobility-derived topology. In a conventional popularity baseline, each edge node either stores its own popular objects or every node stores the same globally popular objects. Those rules are simple, but they ignore an important property of edge networks: nearby nodes can cooperate. If a user frequently moves between two APs, or if two APs are connected through a short handoff path, content stored at one AP can reduce origin access for requests arriving at the other AP. The caching decision is therefore not only a question of what is popular; it is also a question of where the content is stored relative to likely mobility paths.

The study addresses this goal through four linked components. First, mobility traces are converted into a weighted AP graph so that cooperation is represented explicitly. Second, a stochastic objective measures access cost, replica cost, redundancy, and relocation cost under perturbed topology. Third, several policy families are implemented so that the proposed method is compared against meaningful alternatives rather than weak baselines only. Fourth, an adaptive selector chooses among strong candidate policies when topology perturbation changes the operating regime. The main research question is not simply whether a DQN can learn a cache placement. The larger question is whether graph-aware optimization, learning-based refinement, and adaptive selection together produce a coherent caching framework for changing IoT edge networks.

\subsection{How Is the Algorithm Adapting?}

Adaptation occurs at three levels. The first level is demand adaptation. Local and global popularity distributions influence candidate content choices, and the greedy stage evaluates whether each node-content insertion reduces the objective. A content object with high demand is likely to be useful, but it is not automatically selected everywhere because redundancy and cooperative reachability also matter.

The second level is topology adaptation. The greedy policy estimates marginal gains over sampled perturbations of the graph. A placement is rewarded when it serves demand through short cooperative paths and penalized when it duplicates content at strongly related nodes without enough access-cost benefit. This means the same demand matrix can lead to different placements when the graph changes. A central node may be valuable in a stable graph, while a broader distribution of objects may be safer when the topology is moderately uncertain.

The third level is regime adaptation. Adaptive TACC does not assume that one placement family dominates all conditions. It evaluates candidate policies on held-out perturbation samples and chooses the policy with the best validation score after accounting for mean objective, objective variance, and high-churn replication cost. In the reported results, this produces a meaningful switch: DQN-refined greedy is selected at \(p=0.00\) and \(p=0.15\), while diversified popularity is selected at \(p=0.05\) and \(p=0.10\). The algorithm is therefore adapting not by continuously rewriting every cache entry, but by selecting the placement logic that best matches the current perturbation regime.

\subsection{What Changes When the Graph Is Perturbed?}

Graph perturbation changes the active cooperative system. When nodes are removed, the affected AP caches no longer serve requests in that evaluation sample. Any content stored only at removed nodes becomes unavailable to active nodes. When edges are removed, shortest paths may become longer or disappear, which increases cooperative retrieval latency and can convert a cooperative hit into an origin miss. Since the largest connected component is retained, the evaluation focuses on the surviving cooperative region rather than treating isolated fragments as equally usable.

These changes affect policies differently. A placement that concentrates useful content at central APs can perform well in the nominal graph, but it becomes vulnerable if those APs or their incident edges are perturbed. A diversified placement may have slightly worse cost on the stable graph because it is less precisely optimized, yet it can maintain coverage when moderate perturbation removes some paths. A DQN-refined greedy placement can be strong when its local modifications improve access cost without excessive relocation. The adaptive selector exists because these effects are not monotonic: increasing perturbation does not simply make every graph-aware policy worse by the same amount.

\subsection{Why the Adaptive Result Makes Sense}

The measured selector behavior is consistent with the objective. At \(p=0.00\), the graph is stable, so precise topology-aware placement is valuable. DQN refinement improves the greedy warm start by finding local swaps that reduce access cost while keeping relocation small. At \(p=0.05\) and \(p=0.10\), moderate topology changes reduce the reliability of exact central placements. Diversified popularity improves cooperative catalog coverage, so active nodes are more likely to find requested objects somewhere in the surviving graph. At \(p=0.15\), the high-churn penalty discourages fully populated diversified placement because replication overhead becomes more important. The selector therefore returns to DQN-refined greedy.

This pattern is important because it avoids an overclaim. A weaker study could present only the nominal setting and conclude that learning-based refinement is universally best. The perturbation results show a more careful conclusion: learning improves a strong topology-aware warm start in some regimes, diversity is competitive in others, and a practical system should select among candidate policies based on operating conditions. This is the main reason Adaptive TACC is framed as a selector over policy families rather than as a single monolithic neural policy.

\subsection{Sanity Checks for the Results}

Several sanity checks support the interpretation. Global popularity is expected to have low cooperative diversity because every node stores the same high-popularity objects. The results match this expectation: global popularity has poor hit ratio and high objective. Local popularity is expected to produce local hits but miss cooperative opportunities; this is also observed. Random placement sometimes beats local popularity in a dense graph because random diversity can increase the chance that some object is reachable through cooperation, but random placement remains much worse than policies that use demand and topology. Topology-aware greedy improves substantially over random and popularity baselines because it directly optimizes the multi-term objective. DQN refinement gives the largest improvement when the greedy placement is already close enough that local toggles can help.

The relocation metric provides another sanity check. If DQN achieved lower objective only by rewriting many cache entries, the improvement would be operationally less convincing. The nominal DQN result changes only a small number of objects per active node while improving objective and hit ratio. This supports the interpretation that learning is acting as a bounded local refiner rather than an unconstrained cache-rewriting mechanism.

\subsection{Reproducibility Protocol}

The reproducibility protocol is intentionally explicit. The data stage constructs a mobility graph from available movement data or deterministic trace-compatible mobility data. The demand stage uses a fixed Zipf parameter and locality multiplier. The policy stage computes all baselines, greedy placement, DQN refinement, and adaptive selection under the same graph and demand inputs. The evaluation stage repeats each policy over multiple perturbation samples for each perturbation rate. The visualization stage reads the final CSV metrics and generates LaTeX tables and TikZ figures that are included directly in the manuscript.

This protocol reduces the risk that the paper and code disagree. Tables such as the compact metrics table and figures such as the objective trend, hit breakdown, cost components, mobility graph, and selector pie are generated from the same result files used in the analysis. If future experiments change the seed, graph size, content catalog, DQN episodes, or perturbation rates, rerunning the visualization script updates the paper-ready artifacts. The paper therefore functions as both a research report and a reproducibility record.
